\subsection{Predeclared low-column-fraction screening pilot}
\label{sec:lowfraction-round3}
A separate pilot tests whether lower-fraction training repairs the material correction; the primary reconstruction and production coefficients remain unchanged. Twelve training rays use imposed uniform fractions $0.025$, $0.075$ and $0.15$, with 4096 histories per ray/state. A two-coefficient candidate is frozen before testing 12 disjoint rays through four predeclared uniform or heterogeneous fields. Each of the resulting 48 cases has eight paired replicates, with 1024 histories per dry, wet and perturbed state. All 1769472 training and validation histories use the fixed straight-path model with scattering disabled. These are different held-out cases from the earlier $24.98\%\to12.94\%$ comparison; those percentages cannot be used as this pilot's baseline.

The column variable $q$ is water mass divided by total rock--water column mass, including remaining dry rock and excluding air. Training covers $q\in[0.025,0.15]$; held-out values span $[0.0228050475,0.1218551078]$, with one of 48 cases below the training minimum and none above. Imposed local fractions and column fractions are distinct. The dry correction, primary altitude of 30\,km, quadrature and historical control are fixed. Since the material multiplier is unity at $q=0$, this candidate cannot change the dry prediction or repair a dry discrepancy.

Three targets are kept separate. The material-ratio residual compares CSDA wet/dry flux with the ratio of mean MC wet/dry fluxes; the absolute-flux residual compares CSDA wet flux with mean MC wet flux. The response bias is the mean difference between matched-energy CSDA and MC finite-contrast secants in one geometrically preselected crossed cell, at width $\Delta w=0.02$. It has units of integrated directional flux, $\mathrm{m}^{-2}\mathrm{s}^{-1}\mathrm{sr}^{-1}$, per unit local water mass fraction. This is an \emph{absolute} response-bias metric, distinct from relative derivative metrics reported for previous validation sets. Local AD--FD checks of the deterministic midpoint operator are separate from this finite-contrast stochastic target.

Whole-replicate deletion preserves paired-arm, wet/dry and cross-case covariance in the descriptive jackknife SE. Figure~\ref{fig:lowfraction-round3} shows candidate-minus-control changes in absolute bias, with twice these SEs; they are not confidence intervals or calibrated coverage statements. Median absolute material-ratio residual changes only from $2.73133\%$ to $2.71873\%$. The relative median gain is $0.46145\%$, with descriptive SE $29.0589$ percentage points, so the predeclared $30\%$ gain requirement is unresolved. Median absolute-flux residual worsens from $10.5956\%$ to $11.4585\%$, and median absolute response bias from $4.84937\times10^{-5}$ to $5.05680\times10^{-5}$ in the units above. Under the locked screen, 42 of 48 absolute-flux cases and 13 of 48 response cases worsen by more than twice the paired-change SE. These failures determine the screening outcome: \textbf{reject the candidate}; no adoption or held-out retuning follows.

Ninety of 96 arm/case local derivative ladders satisfy the predeclared consistency check. Six larger-step crossing candidates remain unresolved despite agreement at the smallest step; this is not evidence of incorrect AD and is not the substantive rejection cause. Independent raw/provenance checks and a separate implementation reproduce the locked analysis. Two preselected cases replay all 8192 original-seed dry/wet/low/high histories bitwise, and both deterministic arms reproduce; this verification adds no independent replicate.

The accompanying dry diagnostic changes the previous worst ray by only $-0.07932\%$ between 32 and 256 energy nodes, with maximum absolute change $0.09051\%$ over 18 rays. These observed quadrature changes do not account for the historical $59.936\%$ dry discrepancy; 256 nodes are not truth or an asymptotic guarantee. That worst ray also lies outside the dry calibration's angle and slant support, but geometric extrapolation alone does not prove causation. Source-boundary adequacy at 30\,km remains unresolved. This rejected screening pilot leaves the primary correction and reconstruction unchanged, does not establish water recovery, and does not resolve the stationarity limitations in Section~\ref{sec:convergence-round3}.
\begin{figure}[t]\centering
\includegraphics[width=\textwidth]{revision/lowfraction_round3_20261025/fig_lowfraction_round3.pdf}
\caption{Predeclared low-$q$ candidate versus the unchanged historical control on 48 new held-out cases. Each point is the candidate-minus-control change in absolute bias: (a) wet/dry material ratio, (b) absolute wet flux, and (c) the matched-energy finite-contrast response. Ratio and flux changes are percentage points; response changes use absolute flux units per unit local mass fraction. Positive values indicate worsening. Bars are twice the descriptive paired whole-replicate jackknife SE, not confidence limits; colours distinguish the four field configurations. The shaded training column-fraction interval is $[0.025,0.15]$; the single below-range case is retained. The locked screen rejects the candidate, and production coefficients are not changed.}
\label{fig:lowfraction-round3}\end{figure}
